\section{分类训练和限时训练}

\subsection{使用说明}
\begin{enumerate}
  \item \textbf{分类训练}：按知识点逐项突破，标注难度（★基础 / ★★中档 / ★★★拔高）
  \item \textbf{限时训练}：严格控制时间，模拟考场节奏
  \item 先独立完成后对照解析（解析见本书配套答案册）
\end{enumerate}

\subsection{一、分子与细胞}

\begin{enumerate}
  \item[★] 以下哪项不是组成蛋白质的氨基酸？
  \begin{enumerate}
    \item[A.] 含一个氨基和一个羧基的有机物
    \item[B.] 氨基和羧基连在同一个碳原子上
    \item[C.] 含有 C、H、O、N 四种元素
    \item[D.] 结构简式为 R--CH(NH$_2$)--COOH
  \end{enumerate}
  \vspace{0.8cm}

  \item[★] 线粒体和叶绿体的共同点不包括
  \begin{enumerate}
    \item[A.] 均含少量 DNA
    \item[B.] 均含有光合色素
    \item[C.] 均具有双层膜
    \item[D.] 均能产生 ATP
  \end{enumerate}
  \vspace{0.8cm}

  \item[★] 下列关于被动运输的叙述，正确的是
  \begin{enumerate}
    \item[A.] 被动运输需要 ATP 供能
    \item[B.] 自由扩散和协助扩散均顺浓度梯度
    \item[C.] 协助扩散不需要载体蛋白
    \item[D.] O$_2$ 进入细胞的方式为协助扩散
  \end{enumerate}
  \vspace{0.8cm}

  \item[★★] 关于酶的叙述，错误的是
  \begin{enumerate}
    \item[A.] 酶的基本单位是氨基酸或核糖核苷酸
    \item[B.] 酶的作用机理是降低化学反应的活化能
    \item[C.] 温度过高或过低均使酶变性失活
    \item[D.] 酶在细胞内和细胞外均可发挥作用
  \end{enumerate}
  \vspace{0.8cm}

  \item[★★] 将某植物细胞置于 0.3 g/mL 蔗糖溶液中，细胞发生质壁分离。此过程中
  \begin{enumerate}
    \item[A.] 细胞液浓度逐渐升高
    \item[B.] 细胞壁收缩程度大于原生质层
    \item[C.] 细胞失水速率逐渐增大
    \item[D.] 水分从蔗糖溶液向细胞内流动
  \end{enumerate}
  \vspace{0.8cm}

  \item[★★] 右图表示某植物光合速率与光照强度的关系。下列叙述正确的是
  \begin{enumerate}
    \item[A.] a 点表示植物在该光照强度下光合速率最大
    \item[B.] b 点时光合速率等于呼吸速率
    \item[C.] c 点后限制因素主要是光照强度
    \item[D.] d 点时光合作用不再进行
  \end{enumerate}
  \vspace{0.8cm}

  \item[★★] 有氧呼吸与无氧呼吸的共同点不包括
  \begin{enumerate}
    \item[A.] 第一阶段均在细胞质基质进行
    \item[B.] 均能产生少量 ATP
    \item[C.] 均能产生 CO$_2$
    \item[D.] 均需要酶的催化
  \end{enumerate}
  \vspace{0.8cm}

  \item[★★★] 将一株绿色植物置于密闭透明玻璃罩内，在适宜条件下培养。若突然停止 CO$_2$ 供应，短时间内叶绿体中 C$_3$ 和 C$_5$ 含量的变化是
  \begin{enumerate}
    \item[A.] C$_3$ 增加，C$_5$ 减少
    \item[B.] C$_3$ 减少，C$_5$ 增加
    \item[C.] C$_3$ 和 C$_5$ 均减少
    \item[D.] C$_3$ 和 C$_5$ 均增加
  \end{enumerate}
  \vspace{1cm}
\end{enumerate}

\subsection{二、遗传与进化}

\begin{enumerate}
  \item[★] 下列关于 DNA 分子结构的叙述，错误的是
  \begin{enumerate}
    \item[A.] DNA 由两条反向平行的链组成
    \item[B.] 外侧是磷酸和脱氧核糖交替排列
    \item[C.] 碱基之间通过肽键连接
    \item[D.] A 与 T 之间形成两个氢键
  \end{enumerate}
  \vspace{0.8cm}

  \item[★] 基因分离定律的实质是
  \begin{enumerate}
    \item[A.] 子二代出现 3:1 的性状分离比
    \item[B.] 测交后代出现 1:1 的比例
    \item[C.] 等位基因随同源染色体的分离而分开
    \item[D.] 非等位基因自由组合
  \end{enumerate}
  \vspace{0.8cm}

  \item[★★] 豌豆的高茎（D）对矮茎（d）为显性，黄色（Y）对绿色（y）为显性，且独立遗传。基因型为 DdYy 的植株自交，后代中高茎绿色植株占
  \begin{enumerate}
    \item[A.] 9/16
    \item[B.] 3/16
    \item[C.] 1/4
    \item[D.] 1/16
  \end{enumerate}
  \vspace{0.8cm}

  \item[★★] 下列关于基因突变的叙述，正确的是
  \begin{enumerate}
    \item[A.] 基因突变一定会导致生物性状改变
    \item[B.] 基因突变是定向的
    \item[C.] 基因突变只发生在细胞分裂的间期
    \item[D.] DNA 碱基对的替换、增添和缺失均属于基因突变
  \end{enumerate}
  \vspace{0.8cm}

  \item[★★] 红绿色盲为伴 X 隐性遗传病。一对表现型正常的夫妇生了一个色盲男孩，则这对夫妇再生一个色盲女孩的概率是
  \begin{enumerate}
    \item[A.] 0
    \item[B.] 1/4
    \item[C.] 1/2
    \item[D.] 3/4
  \end{enumerate}
  \vspace{0.8cm}

  \item[★★] 下列关于现代生物进化理论的叙述，错误的是
  \begin{enumerate}
    \item[A.] 种群是生物进化的基本单位
    \item[B.] 自然选择决定进化的方向
    \item[C.] 突变和基因重组为进化提供原材料
    \item[D.] 隔离是物种形成的必要条件，生殖隔离一旦形成，个体间不能自由交配
  \end{enumerate}
  \vspace{0.8cm}

  \item[★★★] 某种植物的花色由两对独立遗传的基因控制，存在显性基因时表现为红花，否则为白花。两株红花植株杂交，F$_1$ 中红花:白花 = 15:1，则亲本的基因型为
  \begin{enumerate}
    \item[A.] AABB $\times$ aabb
    \item[B.] AaBb $\times$ AaBb
    \item[C.] AaBB $\times$ Aabb
    \item[D.] AABb $\times$ AaBB
  \end{enumerate}
  \vspace{1cm}
\end{enumerate}

\subsection{三、稳态与调节}

\begin{enumerate}
  \item[★] 内环境稳态的调节机制是
  \begin{enumerate}
    \item[A.] 神经调节
    \item[B.] 体液调节
    \item[C.] 免疫调节
    \item[D.] 神经--体液--免疫调节网络
  \end{enumerate}
  \vspace{0.8cm}

  \item[★] 兴奋在神经纤维上传导时，动作电位的产生是由于
  \begin{enumerate}
    \item[A.] K$^+$ 外流
    \item[B.] Na$^+$ 内流
    \item[C.] Cl$^-$ 内流
    \item[D.] Ca$^{2+}$ 内流
  \end{enumerate}
  \vspace{0.8cm}

  \item[★★] 下列关于人体激素调节的叙述，正确的是
  \begin{enumerate}
    \item[A.] 激素通过导管进入血液
    \item[B.] 激素作用的靶细胞都有相应受体
    \item[C.] 激素直接参与细胞代谢
    \item[D.] 激素在体内含量多、作用高效
  \end{enumerate}
  \vspace{0.8cm}

  \item[★★] 当人体血糖浓度升高时，下列叙述正确的是
  \begin{enumerate}
    \item[A.] 胰岛 A 细胞分泌胰岛素增加
    \item[B.] 肝糖原合成减少
    \item[C.] 组织细胞摄取葡萄糖速率加快
    \item[D.] 胰高血糖素分泌增加
  \end{enumerate}
  \vspace{0.8cm}

  \item[★★] 下列关于免疫的叙述，错误的是
  \begin{enumerate}
    \item[A.] 体液免疫中，浆细胞产生抗体
    \item[B.] 细胞免疫中，细胞毒性 T 细胞裂解靶细胞
    \item[C.] 辅助 T 细胞在特异性免疫中发挥关键作用
    \item[D.] 记忆细胞不能识别抗原
  \end{enumerate}
  \vspace{0.8cm}

  \item[★★] 下列关于生长素的叙述，正确的是
  \begin{enumerate}
    \item[A.] 生长素的化学本质是蛋白质
    \item[B.] 生长素的运输方向是形态学下端到上端
    \item[C.] 生长素的作用具有两重性
    \item[D.] 生长素只能从植物体内合成
  \end{enumerate}
  \vspace{0.8cm}

  \item[★★★] 人在剧烈运动后，肌肉酸疼的主要原因是
  \begin{enumerate}
    \item[A.] 有氧呼吸产生酒精
    \item[B.] 无氧呼吸产生乳酸
    \item[C.] 有氧呼吸产生 CO$_2$
    \item[D.] 无氧呼吸产生 ATP 不足
  \end{enumerate}
  \vspace{1cm}
\end{enumerate}

\subsection{四、生物与环境}

\begin{enumerate}
  \item[★] 下列不属于种群特征的是
  \begin{enumerate}
    \item[A.] 出生率
    \item[B.] 物种丰富度
    \item[C.] 年龄结构
    \item[D.] 性别比例
  \end{enumerate}
  \vspace{0.8cm}

  \item[★] 在食物链"草 $\to$ 兔 $\to$ 狼"中，兔属于
  \begin{enumerate}
    \item[A.] 生产者
    \item[B.] 初级消费者
    \item[C.] 次级消费者
    \item[D.] 分解者
  \end{enumerate}
  \vspace{0.8cm}

  \item[★★] 下列关于生态系统能量流动的叙述，错误的是
  \begin{enumerate}
    \item[A.] 能量沿食物链单向流动
    \item[B.] 能量流动逐级递减，传递效率约为 10--20\%
    \item[C.] 生产者的能量全部流入初级消费者
    \item[D.] 能量流动和物质循环相伴而行
  \end{enumerate}
  \vspace{0.8cm}

  \item[★★] 群落的垂直结构主要体现了
  \begin{enumerate}
    \item[A.] 种群密度的大小
    \item[B.] 物种的丰富度
    \item[C.] 生物对光照等资源的利用
    \item[D.] 种间关系的复杂程度
  \end{enumerate}
  \vspace{0.8cm}

  \item[★★] 下列关于碳循环的叙述，错误的是
  \begin{enumerate}
    \item[A.] 碳在无机环境中主要以 CO$_2$ 形式存在
    \item[B.] 碳从无机环境进入生物群落主要靠光合作用
    \item[C.] 碳在生物群落内部以有机物形式流动
    \item[D.] 碳循环在生物群落和无机环境之间是单向的
  \end{enumerate}
  \vspace{1cm}
\end{enumerate}

\subsection{五、生物技术与工程}

\begin{enumerate}
  \item[★] 制作果酒时，酒精发酵阶段应控制的发酵条件是
  \begin{enumerate}
    \item[A.] 持续通入无菌空气
    \item[B.] 密封、18--25℃
    \item[C.] 加入大量食盐
    \item[D.] 高温灭菌
  \end{enumerate}
  \vspace{0.8cm}

  \item[★] 下列关于 PCR 技术的叙述，错误的是
  \begin{enumerate}
    \item[A.] 需要 DNA 模板
    \item[B.] 需要一对引物
    \item[C.] 需要 DNA 连接酶
    \item[D.] 需要 dNTP 作为原料
  \end{enumerate}
  \vspace{0.8cm}

  \item[★★] 基因工程中，将目的基因导入植物细胞常用的方法是
  \begin{enumerate}
    \item[A.] Ca$^{2+}$ 处理法
    \item[B.] 农杆菌转化法
    \item[C.] 显微注射法
    \item[D.] 脂质体法
  \end{enumerate}
  \vspace{0.8cm}

  \item[★★] 下列关于单克隆抗体的制备，正确的是
  \begin{enumerate}
    \item[A.] 需将抗原注射到正常小鼠体内
    \item[B.] 需将 B 细胞与骨髓瘤细胞融合
    \item[C.] 只需在 HAT 培养基上筛选一次
    \item[D.] 单克隆抗体具有特异性弱、产量高的特点
  \end{enumerate}
  \vspace{0.8cm}

  \item[★★] 植物组织培养中，脱分化是指
  \begin{enumerate}
    \item[A.] 已分化的细胞恢复分裂能力形成愈伤组织
    \item[B.] 愈伤组织形成根和芽
    \item[C.] 外植体直接形成完整植株
    \item[D.] 细胞通过减数分裂形成配子
  \end{enumerate}
  \vspace{1cm}
\end{enumerate}

\subsection{六、限时训练}

\subsubsection{限时训练一（基础卷）}

\textbf{时间}：25 分钟 \quad \textbf{分值}：40 分

\textbf{选择题（每题 2.5 分，共 25 分）}

\begin{enumerate}
  \item 下列物质中，属于还原糖的是
  \begin{enumerate}
    \item[A.] 葡萄糖
    \item[B.] 蔗糖
    \item[C.] 淀粉
    \item[D.] 纤维素
  \end{enumerate}

  \item 有丝分裂后期，细胞中染色体的主要变化是
  \begin{enumerate}
    \item[A.] 染色体复制
    \item[B.] 着丝粒分裂，姐妹染色单体分开
    \item[C.] 同源染色体分离
    \item[D.] 染色体变成染色质
  \end{enumerate}

  \item 下列关于 RNA 的叙述，错误的是
  \begin{enumerate}
    \item[A.] RNA 一般由一条链组成
    \item[B.] RNA 含有 A、G、C、U 四种碱基
    \item[C.] mRNA 是翻译的模板
    \item[D.] 所有 RNA 都参与蛋白质合成
  \end{enumerate}

  \item 当人处于寒冷环境时，机体的调节方式是
  \begin{enumerate}
    \item[A.] 仅神经调节
    \item[B.] 仅体液调节
    \item[C.] 神经--体液调节
    \item[D.] 免疫调节
  \end{enumerate}

  \item 下列属于分解者的是
  \begin{enumerate}
    \item[A.] 小麦
    \item[B.] 蘑菇
    \item[C.] 老虎
    \item[D.] 蝗虫
  \end{enumerate}

  \item 基因型为 Aa 的个体自交，后代纯合子的比例为
  \begin{enumerate}
    \item[A.] 1/2
    \item[B.] 1/4
    \item[C.] 1/3
    \item[D.] 1/8
  \end{enumerate}

  \item 植物细胞发生质壁分离的内因是
  \begin{enumerate}
    \item[A.] 细胞壁的伸缩性大于原生质层
    \item[B.] 原生质层的伸缩性大于细胞壁
    \item[C.] 外界溶液浓度大于细胞液浓度
    \item[D.] 细胞液浓度大于外界溶液浓度
  \end{enumerate}

  \item 下列关于叶绿体色素的叙述，错误的是
  \begin{enumerate}
    \item[A.] 叶绿素 a 呈蓝绿色
    \item[B.] 类胡萝卜素主要吸收红光
    \item[C.] 色素提取可用无水乙醇
    \item[D.] 纸层析法可分离四种色素
  \end{enumerate}

  \item 突触小泡释放的神经递质
  \begin{enumerate}
    \item[A.] 通过主动运输穿过突触间隙
    \item[B.] 与突触后膜上的受体结合
    \item[C.] 只能引起突触后膜兴奋
    \item[D.] 发挥作用后被回收或降解
  \end{enumerate}

  \item 在 DNA 分子中，若 A 占 20\%，则 G 占
  \begin{enumerate}
    \item[A.] 20\%
    \item[B.] 30\%
    \item[C.] 40\%
    \item[D.] 60\%
  \end{enumerate}
\end{enumerate}

\textbf{非选择题（共 15 分）}

11.（7 分）绿色植物光合作用过程可分为光反应和暗反应两个阶段。请回答：
（1）光反应的场所是\_\_\_\_\_\_，产物有\_\_\_\_\_\_、\_\_\_\_\_\_ 和 O$_2$。
（2）暗反应中 CO$_2$ 的固定是指\_\_\_\_\_\_ 的过程。
（3）若突然降低 CO$_2$ 浓度，C$_3$ 含量\_\_\_\_\_\_，C$_5$ 含量\_\_\_\_\_\_。
\vspace{2cm}

12.（8 分）某植物花色由一对等位基因 A/a 控制，红花对白花为显性。现有纯合红花与白花植株杂交，F$_1$ 全为红花，F$_1$ 自交得 F$_2$。
（1）F$_1$ 的基因型为\_\_\_\_\_\_。
（2）F$_2$ 的表现型及比例为\_\_\_\_\_\_。
（3）F$_2$ 中红花植株中纯合子占\_\_\_\_\_\_。
（4）欲鉴定一株红花植株是否为纯合子，最简便的方法是\_\_\_\_\_\_。
\vspace{2cm}

\subsubsection{限时训练二（拔高卷）}

\textbf{时间}：30 分钟 \quad \textbf{分值}：50 分

\textbf{选择题（每题 2.5 分，共 20 分）}

\begin{enumerate}
  \item 下列关于细胞生命历程的叙述，正确的是
  \begin{enumerate}
    \item[A.] 细胞分化过程中遗传物质发生改变
    \item[B.] 细胞癌变后细胞周期变长
    \item[C.] 细胞衰老时酶活性普遍降低
    \item[D.] 细胞凋亡是被动死亡过程
  \end{enumerate}

  \item 某种群中 AA 占 20\%，Aa 占 60\%，aa 占 20\%，则该种群中 A 的基因频率为
  \begin{enumerate}
    \item[A.] 40\%
    \item[B.] 50\%
    \item[C.] 60\%
    \item[D.] 80\%
  \end{enumerate}

  \item 下列关于生态系统信息传递的叙述，错误的是
  \begin{enumerate}
    \item[A.] 物理信息可以来自非生物环境
    \item[B.] 化学信息通过信息素传递
    \item[C.] 行为信息只发生在同种生物之间
    \item[D.] 信息传递对生态系统稳定性有重要作用
  \end{enumerate}

  \item 某双链 DNA 分子含 1000 个碱基对，其中 A 占 30\%，则该 DNA 分子中 G 的数目为
  \begin{enumerate}
    \item[A.] 200
    \item[B.] 400
    \item[C.] 600
    \item[D.] 800
  \end{enumerate}

  \item 下列关于植物激素的叙述，错误的是
  \begin{enumerate}
    \item[A.] 赤霉素能促进种子萌发
    \item[B.] 脱落酸能抑制细胞分裂
    \item[C.] 乙烯能促进果实成熟
    \item[D.] 细胞分裂素能促进叶的衰老脱落
  \end{enumerate}

  \item 下列关于进化的叙述，正确的是
  \begin{enumerate}
    \item[A.] 自然选择决定了生物变异的方向
    \item[B.] 基因频率改变一定产生新物种
    \item[C.] 隔离是物种形成的必要条件
    \item[D.] 突变和基因重组决定进化方向
  \end{enumerate}

  \item 在啤酒酿造过程中，利用到的微生物主要是
  \begin{enumerate}
    \item[A.] 醋酸菌
    \item[B.] 酵母菌
    \item[C.] 乳酸菌
    \item[D.] 毛霉
  \end{enumerate}

  \item 下列关于动物细胞核移植技术的叙述，正确的是
  \begin{enumerate}
    \item[A.] 供体细胞应选用体细胞
    \item[B.] 受体细胞应选用受精卵
    \item[C.] 核移植技术的原理是细胞全能性
    \item[D.] 克隆动物与供体动物性状完全相同
  \end{enumerate}
\end{enumerate}

\textbf{非选择题（共 30 分）}

9.（10 分）人体在剧烈运动时，机体通过调节维持稳态。请回答：
（1）剧烈运动时，细胞主要通过\_\_\_\_\_\_（方式）产生能量，产生的 CO$_2$ 会刺激位于\_\_\_\_\_\_的呼吸中枢，使呼吸加深加快。
（2）血糖浓度降低时，胰岛\_\_\_\_\_\_细胞分泌胰高血糖素增加，促进\_\_\_\_\_\_ 分解为葡萄糖。
（3）大量出汗后，细胞外液渗透压\_\_\_\_\_\_，刺激下丘脑渗透压感受器，引起\_\_\_\_\_\_激素释放增加，促进肾小管和集合管对水的重吸收。
\vspace{3cm}

10.（10 分）某生态系统中存在食物链：草 $\to$ 蝗虫 $\to$ 食虫鸟 $\to$ 鹰。回答：
（1）该食物链中，食虫鸟属于第\_\_\_\_\_\_营养级，鹰属于\_\_\_\_\_\_消费者。
（2）若鹰增加 1 kg 体重，至少需要消耗草\_\_\_\_\_\_kg（传递效率按 20\% 计）。
（3）信息传递在生态系统中的作用包括：\_\_\_\_\_\_、\_\_\_\_\_\_、\_\_\_\_\_\_。
\vspace{3cm}

11.（10 分）现有纯合高茎抗病和矮茎易感病两种小麦（两对基因独立遗传，高茎 D 对矮茎 d 为显性，抗病 R 对易感病 r 为显性）。请设计育种方案，获得矮茎抗病纯合品种。

\textbf{方案}：
\begin{enumerate}
  \item \_\_\_\_\_\_
  \item \_\_\_\_\_\_
  \item \_\_\_\_\_\_
  \item \_\_\_\_\_\_
\end{enumerate}
